% IEEE CloudNet 2026 — Demo & Artifact Track
% 2 pages, 10 pt, double-column
% Compile: pdflatex main && pdflatex main
%
\documentclass[10pt,conference]{IEEEtran}

\usepackage[T1]{fontenc}
\usepackage[utf8]{inputenc}
\usepackage{booktabs}
\usepackage{url}
\usepackage{cite}
\usepackage{makecell}
\usepackage[table]{xcolor}

% ─────────────────────────────────────────────────────────────────────────────
\title{QoS-Aware AI Inference Offloading for Vehicular Edge Computing:\\
       A Physics-Based Dataset and Simulation Framework}

\author{%
  \IEEEauthorblockN{%
    Konstantinos Gyftodimos,\quad
    Eirini Liotou,\quad
    Kyriakos Chiotis,\quad
    George Dimitrakopoulos%
  }%
  \IEEEauthorblockA{%
    Dept.\ of Informatics and Telematics,
    Harokopio University of Athens (HUA), Athens, Greece\\
    \{kgyftodimos,\;eliotou,\;kchiotis,\;gdimitra\}@\textit{hua.gr}\quad
    \textit{Presenter: K.\,Gyftodimos\,---\,onsite}%
  }%
}

% ─────────────────────────────────────────────────────────────────────────────
\begin{document}
\maketitle

% ── Abstract ─────────────────────────────────────────────────────────────────
\begin{abstract}
We present a physics-based open-source simulation framework for QoS-aware AI
inference offloading in vehicular edge computing, comprising a 50,000-sample
synthetic dataset across five road scenarios, a multi-output machine-learning
predictor, and a real-time closed-loop demonstration.
\end{abstract}

% ── I. Overview ──────────────────────────────────────────────────────────────
\section{Overview}

Autonomous vehicles generate continuous high-rate video streams that must be
processed by YOLOv8~\cite{yolov8} object detection under a strict 100\,ms
end-to-end latency constraint.  The simulation framework models a four-tier
execution hierarchy: (A)~a dedicated lightweight on-board neural processing
unit (NPU), (B)~the vehicle's primary high-throughput GPU, (C)~a nearby
Multi-Access Edge Computing (MEC)~\cite{etsi_mec} server, and (D)~a remote
cloud data centre.  Each tier runs a different YOLOv8 variant and offers a
distinct latency--accuracy trade-off (Table~\ref{tab:modes}).  Selecting the
optimal tier in real time requires jointly predicting network conditions and
enforcing Quality of Service (QoS) constraints simultaneously.

Most vehicular offloading studies rely on proprietary radio traces or abstract
the channel to the point that MEC latency variability is not faithfully
reproduced.  This artifact provides an open, end-to-end reproducible
baseline---from 3GPP-calibrated channel simulation to closed-loop
optimisation---that is fully self-contained, runs offline, and is portable to
any compute-intensive edge inference workload.  Use cases range from benchmarking novel offloading policies to evaluating
new edge AI models and prototyping QoS-aware placement for autonomous
driving or smart-infrastructure monitoring.

\textbf{Delivery format:} physical, onsite demonstration at IEEE
CloudNet~2026.  The complete system runs as a single terminal command on a
standard laptop with no external hardware or network connection required.

% ── II. Innovation \& Contribution ───────────────────────────────────────────
\section{Innovation \& Contribution}

\begin{enumerate}
  \item \textbf{Physics-based synthetic dataset (50\,k samples).}
        Generated via a 3GPP TR\,38.901-calibrated log-distance path-loss
        model~\cite{3gpp38901} with log-normal shadowing and non-line-of-sight
        (NLOS) transitions, combined with Shannon--Hartley
        throughput~\cite{shannon1948}.  Five road scenarios span 300\,m
        to 10\,km MEC inter-site distances (urban dense through rural),
        with scenario-specific speed, bandwidth, and frequency profiles.
        Fixed-seed (42) reproducible; MIT license; schema in
        Table~\ref{tab:schema}.

  \item \textbf{Multi-output machine-learning (ML) predictor.}
        A \texttt{RandomForestRegressor} (200 trees, 80/20 train-test split)
        predicts MEC and cloud round-trip time (RTT) and server capacity
        utilisation; a \texttt{GradientBoostingClassifier} predicts MEC
        availability (binary).  All scikit-learn~\cite{sklearn} pipelines are
        serialised to disk after training and loaded on subsequent runs,
        reducing cold-start time to $\approx$1\,s.  Per-frame predictions are
        smoothed with an exponential moving average ($\alpha$\,=\,0.3) to
        suppress oscillation under transient channel variations.

  \item \textbf{Edge-first QoS optimizer.}
        The decision engine enforces three hard QoS budgets: total latency
        $\leq$100\,ms, accuracy loss $\leq$15\%, and server capacity
        $<$95\%.  Among feasible modes it applies an edge-first preference
        rule: Mode~C (MEC) is chosen over Mode~D (Cloud) when the
        incremental mean Average Precision (mAP) gain of the cloud model is
        below 5\,pts, thereby conserving Wide Area Network (WAN) bandwidth.
        If no mode satisfies all constraints---e.g.\ during a connectivity
        outage---the optimizer falls back to Mode~B (on-board GPU) with an
        explicit reason flag (\texttt{fallback\_no\_feasible\_mode}).

  \item \textbf{Live four-tier failover demo.}
        The complete pipeline executes in a single command and demonstrates
        all four execution tiers within one 100-frame run.  A simulated
        connectivity outage (frames 40--54) forces MEC unreachable and
        inflates cloud RTT to 200\,ms, driving the system through a full
        C\,$\to$\,B\,$\to$\,C failover and making all QoS constraint paths
        observable in real time.
\end{enumerate}

% Tables II and I placed after Innovation so LaTeX floats both to page 2.
\begin{table}[t]
\centering
\caption{YOLOv8 Execution Mode Specifications}
\label{tab:modes}
\scriptsize\setlength{\tabcolsep}{5pt}
\begin{tabular}{@{}clrrr@{}}
\toprule
\textbf{Mode} & \textbf{Hardware\,/\,Model} &
  \textbf{Tot.\,(ms)} & \textbf{mAP} & \textbf{AccLoss} \\
\midrule
A & Car NPU\,/\,YOLOv8n  &   8       & 37.3 & 20.8\%$^*$ \\
B & Car GPU\,/\,YOLOv8n  &   5       & 37.3 & 20.8\%$^*$ \\
C & MEC\,/\,YOLOv8m      & $\sim$24  & 50.2 &  4.9\%     \\
D & Cloud\,/\,YOLOv8l    & $\sim$84  & 52.9 &  0.0\%     \\
\bottomrule
\multicolumn{5}{@{}l}{Budget: latency $\leq$100\,ms, acc.\ loss $\leq$15\%, cap.\ $<$95\%.}\\
\multicolumn{5}{@{}l}{$^*$Exceed budget; activate only as emergency fallback.}
\end{tabular}
\end{table}

\begin{table}[t]
\centering
\caption{Dataset Feature Schema (key columns; full schema in repository)}
\label{tab:schema}
\small\setlength{\tabcolsep}{4pt}
\begin{tabular}{@{}lll@{}}
\toprule
\textbf{Feature} & \textbf{Type} & \textbf{Description} \\
\midrule
\texttt{scenario}              & cat.  & Road environment (5 types) \\
\texttt{car\_speed\_kmh}       & float & Vehicle speed (km/h) \\
\texttt{mec\_distance\_m}      & float & Distance to MEC tower (m) \\
\texttt{mec\_rssi\_dbm}        & float & Received signal strength (dBm) \\
\texttt{num\_connected\_ues}   & int   & User equipment sharing cell \\
\texttt{cloud\_congestion}     & float & WAN congestion index [0,\,1] \\
\texttt{payload\_kb}           & float & Compressed frame size (kB) \\
\midrule
\texttt{mec\_available}        & bin.  & \textbf{Target:} MEC reachability \\
\texttt{mec\_network\_ms}      & float & \textbf{Target:} MEC RTT (ms) \\
\texttt{mec\_capacity\_pct}    & float & \textbf{Target:} MEC load (\%) \\
\texttt{cloud\_network\_ms}    & float & \textbf{Target:} Cloud RTT (ms) \\
\texttt{cloud\_capacity\_pct}  & float & \textbf{Target:} Cloud load (\%) \\
\bottomrule
\end{tabular}
\end{table}

% ── III. Relevance to CloudNet ────────────────────────────────────────────────
\section{Relevance to CloudNet}

The artifact targets CloudNet researchers working on vehicular edge
computing~\cite{shi2016edge}, task offloading~\cite{chen2019offloading}, and
5G edge AI~\cite{liu2020edge}.  Key takeaways:
\begin{itemize}
  \item \textbf{Reproducible baseline.}  A 3GPP-calibrated MEC offloading
        dataset across five deployment archetypes, from dense urban (300\,m
        MEC spacing) to rural (10\,km), freely available for comparison and
        extension by the research community.
  \item \textbf{Offloading as a necessity.}  Empirical evidence that on-board
        YOLOv8n (Modes~A/B) incurs a 20.8\% accuracy loss relative to cloud
        YOLOv8l, exceeding the 15\% QoS budget and making MEC/cloud
        offloading mandatory whenever connectivity permits.
  \item \textbf{End-to-end blueprint.}  A complete, runnable coupling of ML
        network-state prediction with a policy-based QoS optimizer, portable
        to other inference workloads (e.g.\ LiDAR processing, V2X perception)
        and extendable to new tier definitions or QoS policies.
\end{itemize}

% ── IV. Demonstration Plan \& Requirements ────────────────────────────────────
\section{Demonstration Plan \& Requirements}

\textbf{Setup.}
Install dependencies (\texttt{pip install -r requirements.txt}) and launch
\texttt{python demo.py}.  On first run the system auto-generates the 50\,k
synthetic dataset and trains all ML models ($\approx$50\,s total);
subsequent runs load the serialised models and start in $\approx$2\,s.

\textbf{What the audience observes.}
The terminal table (Table~\ref{tab:live}) refreshes at 5\,Hz for 100 frames,
showing all four tiers simultaneously with live per-mode latency, accuracy
loss, capacity, and decision reason.  Under normal suburban conditions Mode~C
(MEC / YOLOv8m, $\approx$57\,ms total) is selected over Mode~D (Cloud /
YOLOv8l, $\approx$72\,ms) by the edge-first rule.  Modes~A and~B are visible
every frame, marked infeasible due to the 20.8\% accuracy-loss violation.
Starting at frame~40 an injected outage forces MEC unavailable and cloud RTT
to 200\,ms; the optimizer switches to Mode~B (5\,ms on-board GPU) with reason
\texttt{fallback\_no\_feasible\_mode}, demonstrating the full failover path
before recovering to Mode~C once the outage clears.

\textbf{AV / Power:} standard laptop (CPU-only, no GPU required); terminal
with ANSI colour support.  No special power supply, networking, or AV
equipment needed.  Total runtime: $\approx$50\,s first run, $\approx$20\,s
subsequently.

% ── V. Artifact Details ───────────────────────────────────────────────────────
\section{Artifact Details}

All components are in a public GitHub repository
(\url{https://github.com/[org]/cloudnet2026-offloading}, release
\texttt{v1.0.0}; Ubuntu~22.04 / macOS~14 / Python~3.11).

\textbf{Dataset:} 50\,k-row CSV, 13 features\,+\,5 targets
(Table~\ref{tab:schema}); synthetic, MIT license; Zenodo DOI on
acceptance; regenerate via
\texttt{python -m src.dataset.generate -{}-seed 42}.

\textbf{Framework:} \texttt{src/dataset/} (channel sim),
\texttt{src/predictor/} (ML pipelines), \texttt{src/offloading/}
(QoS optimizer).  Dataset and model bundled; auto-recreated if absent.

\textbf{Demo:} \texttt{python demo.py} shows live per-mode latency,
accuracy loss, capacity, and decision reason; configurable via
\texttt{-{}-scenario}, \texttt{-{}-hz}, \texttt{-{}-outage-at}.

% ── VI. Ethics, Privacy \& Safety ────────────────────────────────────────────
\section{Ethics, Privacy \& Safety}

The dataset is entirely synthetic---generated from a calibrated radio-physics
model with no real-world measurements, user data, or personally identifiable
information.  No human subjects or consent procedures are required.  The
framework poses no safety risk.

% ── Acknowledgment ───────────────────────────────────────────────────────────
\section*{Acknowledgment}

This work has been funded by the Shift2SDV project (Grant Agreement
No.~101194245), supported by the Chips Joint Undertaking and its members,
including top-up funding by the national authorities of Austria, Denmark,
Germany, Greece, Finland, Italy, Netherlands, Poland, Portugal, Spain, and
Turkey.

% Table III: single-column, placed before bibliography so LaTeX
% numbers it correctly and floats it to the bottom of a page-2 column.
\begin{table}[b]
\centering
\caption{%
  Live optimizer snapshot (suburban, frame 100/100, MEC dist.\ 611\,m,
  signal $-$55.5\,dBm, 4 mode switches).
  Tot.\,(ms)\,=\,inference\,$+$\,network RTT.%
}
\label{tab:live}
\renewcommand{\arraystretch}{0.95}
\setlength{\tabcolsep}{3pt}
\scriptsize
\begin{tabular}{@{} c l r r r l @{}}
\Xhline{2\arrayrulewidth}
\textbf{M.} & \textbf{Model} &
  \textbf{Tot.\,(ms)} & \textbf{mAP} & \textbf{AccLoss} &
  \textbf{Status} \\
\hline
A & \texttt{yolov8n} &  8.0 & 37.3 & 20.8\% & $\times$\ loss\,$>$\,15\% \\
B & \texttt{yolov8n} &  5.0 & 37.3 & 20.8\% & $\times$\ loss\,$>$\,15\% \\
\rowcolor{gray!15}
C & \texttt{yolov8m} & 56.9 & 50.2 &  4.9\% &
  $\blacktriangleleft$\ \textbf{SELECTED} \\
D & \texttt{yolov8l} & 71.6 & 52.9 &  0.0\% & feasible, not preferred \\
\Xhline{2\arrayrulewidth}
\end{tabular}
\end{table}

\begin{thebibliography}{8}

\bibitem{yolov8}
G.~Jocher, A.~Chaurasia, and J.~Qiu,
``Ultralytics YOLOv8,'' GitHub, 2023.
[Online]. Available: \url{https://github.com/ultralytics/ultralytics}

\bibitem{etsi_mec}
ETSI,
``GS MEC~003: MEC Framework and Reference Architecture, v3.1.1,'' 2022.

\bibitem{3gpp38901}
3GPP,
``TR~38.901: Study on channel model for frequencies from 0.5 to
100~GHz, v16.1.0,'' 3rd Generation Partnership Project, 2019.

\bibitem{shannon1948}
C.~E.~Shannon,
``A mathematical theory of communication,''
\textit{Bell Syst. Tech. J.}, vol.~27, pp.~379--423, 1948.

\bibitem{shi2016edge}
W.~Shi \textit{et~al.},
``Edge computing: Vision and challenges,''
\textit{IEEE Internet Things J.}, vol.~3, no.~5, pp.~637--646, 2016.

\bibitem{chen2019offloading}
X.~Chen \textit{et~al.},
``Optimized computation offloading via deep reinforcement learning,''
\textit{IEEE Internet Things J.}, vol.~6, no.~3, pp.~4005--4018, 2019.

\bibitem{liu2020edge}
Y.~Liu \textit{et~al.},
``Toward edge intelligence: MEC for 5G and IoT,''
\textit{IEEE Internet Things J.}, vol.~7, no.~8, pp.~6722--6747, 2020.

\end{thebibliography}

\end{document}
